\chapter{1961:Melvin Calvin}

\label{chap:chemistry1961}

The Nobel Prize in Chemistry for the year 1961 was awarded to Melvin Calvin.

\textbf{Motivation:}

for his research on the carbon dioxide assimilation in plants

\section{Melvin Calvin}

\subsection*{Early Life and Education}

Melvin Calvin was born on 1911-04-08 in St. Paul, Minnesota, USA.

\subsection*{Career and Major Research}

Calvin studied at the Michigan College of Mining and Technology and received his doctorate from the University of California, Berkeley, in 1935. He spent his career at Berkeley, where he worked at the Lawrence Radiation Laboratory and became professor of chemistry. Using the radioactive tracer carbon-14, introduced into photosynthesis research by his colleague Andrew Benson, Calvin and his group, including James Bassham, traced the path of carbon in photosynthesis and worked out the Calvin cycle.

\subsection*{Nobel Prize-winning Work}

The work for which Melvin Calvin was awarded the Nobel Prize in Chemistry in 1961 was described by the Nobel Committee as: "for his research on the carbon dioxide assimilation in plants".

By exposing algae to carbon-14 labeled carbon dioxide for brief periods and analyzing the products, Calvin's group identified the sequence of reactions by which carbon dioxide is fixed into organic compounds. They showed that the first product is 3-phosphoglycerate and that the cycle regenerates its five-carbon acceptor, ribulose-1,5-bisphosphate, so that carbon can be assimilated continuously.

\subsection*{Significance and Impact}

The Calvin cycle, as it came to be called, was the first complete description of a major metabolic pathway, and it explained how plants convert carbon dioxide into the sugars on which all life depends. The tracer methods used became standard in biochemistry.

\subsection*{Later Career and Honors}

Calvin remained at the University of California, Berkeley, where he directed the Laboratory of Chemical Biodynamics. He died on 1997-01-08 in Berkeley, California.

\subsection*{Legacy}

The Calvin cycle is one of the foundations of biochemistry and plant science, and Calvin's use of isotopic tracers set the pattern for the elucidation of metabolic pathways throughout the biological sciences.

\subsection*{Modern Developments}

Calvin's elucidation of the photosynthetic carbon cycle (the Calvin cycle) revealed how plants convert carbon dioxide into sugar and remains the basis of our understanding of carbon fixation. His work underpins modern efforts to improve crop yields, engineer carbon capture, and develop artificial photosynthesis for renewable fuels. Radioactive tracer techniques used in his research are now standard in biology.

\subsection*{Selected Publications}

\begin{itemize}

\item \emph{The Path of Carbon in Photosynthesis}, Science, 1962.

\end{itemize}

\clearpage

\section*{Chapter Summary}

The 1961 prize recognized research on carbon dioxide assimilation in plants. Melvin Calvin and his colleagues at Berkeley used carbon-14 tracers to trace the path of carbon in photosynthesis and established the cycle of reactions by which plants fix carbon dioxide into carbohydrates.

